% Generated draft: requires mathematical and source review.
\documentclass[11pt]{article}
\usepackage[margin=25mm]{geometry}
\usepackage{fontspec}
\setmainfont{Latin Modern Roman}
\usepackage{amsmath,amssymb,mathtools,unicode-math}
\setmathfont{Latin Modern Math}
\usepackage{enumitem,needspace,xurl,hyperref}
\setlength{\parindent}{0pt}
\setlength{\parskip}{5pt}
\setlength{\emergencystretch}{3em}
\newcommand{\N}{\mathbb{N}}
\newcommand{\Z}{\mathbb{Z}}
\newcommand{\Q}{\mathbb{Q}}
\newcommand{\R}{\mathbb{R}}
\newcommand{\C}{\mathbb{C}}
\newcommand{\F}{\mathbb{F}}
\newcommand{\A}{\mathbb{A}}
\newcommand{\PP}{\mathbb{P}}
\newcommand{\E}{\mathbb{E}}
\newcommand{\eps}{\varepsilon}
\newcommand{\dd}{\,\mathrm{d}}
\begin{document}
\section*{UnsolvedMath ID: 701701}
\textbf{Source problem code:} catalogue entry
\subsection{Definitions and mathematical statement}
Let \(\mathcal G\) be a family of finite groups. For an integer \(r\geq 1\), let \(F_r\) denote the free group on generators \(x_1,\ldots,x_r\). A common group law for \(\mathcal G\) means a nonidentity word \(w\in F_r\) such that
\[
\forall G\in\mathcal G\quad\forall(g_1,\ldots,g_r)\in G^r,\qquad w(g_1,\ldots,g_r)=1_G.
\]
Here the same word \(w\) must work for every group in the family.

Use the following Følner formulation of uniform amenability: for every real \(\varepsilon>0\) and every integer \(n\geq1\), there is an integer \(M\geq1\) such that, for every \(G\in\mathcal G\) and every subset \(S\subseteq G\) with \(|S|\leq n\), there exists a nonempty subset \(A\subseteq G\) satisfying
\[
|A|\leq M,
\qquad
|sA\mathbin{\triangle}A|\leq\varepsilon|A|
\quad\text{for every }s\in S.
\]
Here \(sA=\{sa:a\in A\}\), and \(B\mathbin{\triangle}C=(B\setminus C)\cup(C\setminus B)\). The bound \(M\) may depend on \(\mathcal G\), \(n\), and \(\varepsilon\), but must be independent of \(G\) and \(S\).

\textbf{Question (Kourovka Notebook 20.43).} Is every family \(\mathcal G\) of finite groups admitting a common group law uniformly amenable in this sense?

The supplied record does not explicitly define either a group law or uniform amenability. This draft interprets a group law as a nontrivial word identity and uses the family-wide Følner condition above; these conventions require confirmation against the cited source.
\subsection{Short English statement}
If one fixed nontrivial word is an identity in every finite group of a family, must approximate invariance under any bounded number of translations be achievable by a nonempty set whose size has a bound uniform across the family? The bound may depend on the family, the number of translations, and the allowed error.
\subsection{Sources}
\begin{itemize}
\item Local UnsolvedMath registry: supplied complete source record, permanent ID \(701701\), exact source problem code \texttt{catalogue entry}; upstream revision \texttt{d371ee917ac9647fbb06cd6fc529deec6f3dda82}. The record was created and updated on \texttt{2026-10-16T14:00:00Z}.
\item \emph{The Kourovka Notebook (Unsolved Problems in Group Theory)}, 20th issue (2022), Problem 20.43; also identified in the supplied background with the 21st issue edition (2026), \texttt{arXiv:1401.0300v48}. The supplied collection description names E. I. Khukhro and V. D. Mazurov as editors.
\item Proposed by S. Kionke and E. Schesler, with proposed year 2022 according to the supplied record.
\item The registry labels this problem \emph{open}. This is supplied source metadata, not an independently verified assessment of the literature. No separate literature-check date or literature-source links were supplied.
\end{itemize}
\end{document}
